\begin{frame}{Merkle trees \only<2->{-- Proof}}
    \begin{columns}[c]
        \begin{column}{0.5\textwidth}
            \only<1>{It is a binary tree $T$ of height $H$ with $2^H$ leaves and
            $2^{H} - 1$ internal
nodes. Each leaf stores the cryptographic hash of the underlying data, rather than
the raw data itself. The same cryptographic hash function is applied recursively at
internal nodes, which store the hash of the concatenation of their two children.
            $$n_{parent} = f(n_{left} || n_{right})$$}
\only<2>{
Merkle trees has the ability to prove that a
given piece of data is part of a larger set, without revealing or recomputing the
            entire dataset. For $data_1$ we have $\pi=\{node_0, node_{23}\}$.\\
We can make a proof verification
            comparing the result with the presumed root in $O(\log n)$.
            }
        \end{column}
        \begin{column}{0.5\textwidth}
            \centering
            \begin{tikzpicture}[
              every node/.style={font=\tiny},
              leaf/.style={draw, circle, minimum size=7mm},
              hash/.style={draw, rectangle, rounded corners=2pt,
              minimum width=10mm, minimum height=7mm, 
              red/.style={draw=red},
              sibling distance=1cm,
              align=center},
              node distance=.6cm
            ]
            % leaves
            \node[leaf,] (data0) {$data_0$};
            \only<1>{\node[leaf, right=of data0] (data1) {$data_1$};}
            \only<2->{\node[leaf, red, right=of data0] (data1) {$data_1$};}
            \node[leaf, right=of data1] (data2) {$data_2$};
            \node[leaf, right=of data2] (data3) {$data_3$};
            % leaf-hash nodes with two lines (inline math for each line)
            \only<1>{\node[hash, above=0.2cm of data0] (n0) {\(\mathbf{node_0}\)\\ \(f(data_0)\)};}
            \only<2->{\node[hash, red, above=0.2cm of data0] (n0) {\(\mathbf{node_0}\)\\ \(f(data_0)\)};}
            \node[hash, above=0.2cm of data1] (n1) {\(\mathbf{node_1}\)\\ \(f(data_1)\)};
            \node[hash, above=0.2cm of data2] (n2) {\(\mathbf{node_2}\)\\ \(f(data_2)\)};
            \node[hash, above=0.2cm of data3] (n3) {\(\mathbf{node_3}\)\\ \(f(data_3)\)};
            % parents
            \node[hash, above=of $(n0)!0.5!(n1)$] (p01) {\(\mathbf{node_{01}}\)\\ \(f(node_0\Vert node_1)\)};
            \only<1>{\node[hash, above=of $(n2)!0.5!(n3)$] (p23) {\(\mathbf{node_{23}}\)\\ \(f(node_2\Vert node_3)\)};}
            \only<2->{\node[hash, red, above=of $(n2)!0.5!(n3)$] (p23) {\(\mathbf{node_{23}}\)\\ \(f(node_2\Vert node_3)\)};}
            % root
            \node[hash, above=of $(p01)!0.5!(p23)$] (root) {\(\textbf{root}\)\\ \(f(node_{01}\Vert node_{23})\)};
            % edges with arrows
            \foreach \a/\b in {data0/n0, data1/n1, data2/n2, data3/n3, n0/p01, n1/p01, n2/p23, n3/p23, p01/root, p23/root}
              \draw[->] (\a) -- (\b);
            \end{tikzpicture}
        \end{column}
    \end{columns}
\end{frame}

\begin{frame}{Merkle trees -- Lib}
    \begin{columns}[c]
        \begin{column}{0.7\textwidth}
            A dedicated Rust library was developed to efficiently compute Merkle tree root hashes.
            \begin{table}[h!]
                \centering
                \footnotesize % or \small, \scriptsize, \tiny
                \begin{tabular}{|l|c|c|c|}
                \hline
                    \textbf{Hash function} & \textbf{5 MB} & \textbf{10 MB} & \textbf{15 MB} \\
                \hline
                    SHA-256   & 89.901 ms  & 178.42 ms  & 268.53 ms \\
                    Keccak-256 & 521.49 ms & 1.1334 s   & 1.3438 s \\
                    BLAKE3    & 73.091 ms  & 154.68 ms  & 219.79 ms \\
                \hline
                \end{tabular}
            \end{table}
        \end{column}
        \begin{column}{0.3\textwidth}
            \centering
            \includegraphics[width=.4\linewidth]{static/rust-logo.png}
        \end{column}
    \end{columns}

\end{frame}

\begin{frame}{Raft}
A consensus protocol, where each server on a cluster is a follower, a candidate
or a leader. There is a single leader, responsible for sending messages to other servers via the Raft log.

\centering
\includegraphics[width=.4\linewidth]{static/raft-logo.png}

\end{frame}

\begin{frame}{Reed-Solomon}

File split in $n+k$ pieces. Each piece to a different agent.
\vspace{1em}

\centering

\begin{tikzpicture}[
    % Set a global distance for 'right of' positioning
    node distance=1.5cm 
]

% Node 1: Start at (0,0)
\node[
    draw,
    cylinder,
    cylinder uses custom fill,
    cylinder body fill=gray!10,
    cylinder end fill=gray!40,
    minimum height=2.5cm,
    minimum width=1.5cm,
    align=center
] (C1) at (0,0) {Agent\\ 1};

\node[
    draw,
    cylinder,
    cylinder uses custom fill,
    cylinder body fill=red!70,
    cylinder end fill=red!90,
    minimum height=2.5cm,
    minimum width=1.5cm,
    align=center,
    right=of C1 % <--- Positioning command
] (C2) {Agent\\ 2};

% Node 3: Positioned to the right of Node 2
\node[
    draw,
    cylinder,
    cylinder uses custom fill,
    cylinder body fill=gray!10,
    cylinder end fill=gray!40,
    minimum height=2.5cm,
    minimum width=1.5cm,
    align=center,
    right=of C2 % <--- Positioning command
] (C3) {Agent\\ 3};

\node[
    draw,
    circle,
    fill=blue!30,
    align=center,
    below=of C2 % <--- Positioning command
] (GW) {Gateway};


\draw[->, thick] (GW) -- (C1)
    node[midway, left, font=\scriptsize] {Shard 1};

\draw[->, red, thick] (GW) -- (C2)
    node[midway, left, font=\scriptsize] {Shard 2};

\draw[->, thick] (GW) -- (C3)
    node[midway, right, font=\scriptsize] {Shard 3};

\end{tikzpicture}

\end{frame}
